\documentclass{beamer}
% =====================================================================
% finance-beamer.tex — shared Beamer style for the LLM-in-Finance course
% =====================================================================
% Reusable slide style. Each deck \input's this immediately after
% \documentclass{beamer}. It provides:
%   * the Madrid theme + book colour palette (matches the printed book)
%   * two book-style framing blocks:
%       \begin{bigpicture} ... \end{bigpicture}   ("the bigger picture")
%       \begin{underhood}   ... \end{underhood}    ("under the hood")
%     These mirror the book's `context` and `deepdive` tcolorboxes so the
%     same intuition-vs-mechanism scaffold the reader meets in the book
%     reappears on the slides.
%   * a Python listings style for code frames
%   * appendix conventions: \appendixstart drops a divider and switches
%     to non-numbered backup frames so "extra / less central" material
%     lives after the main talk without inflating the slide count.
%
% Usage:
%   \documentclass{beamer}
%   % =====================================================================
% finance-beamer.tex — shared Beamer style for the LLM-in-Finance course
% =====================================================================
% Reusable slide style. Each deck \input's this immediately after
% \documentclass{beamer}. It provides:
%   * the Madrid theme + book colour palette (matches the printed book)
%   * two book-style framing blocks:
%       \begin{bigpicture} ... \end{bigpicture}   ("the bigger picture")
%       \begin{underhood}   ... \end{underhood}    ("under the hood")
%     These mirror the book's `context` and `deepdive` tcolorboxes so the
%     same intuition-vs-mechanism scaffold the reader meets in the book
%     reappears on the slides.
%   * a Python listings style for code frames
%   * appendix conventions: \appendixstart drops a divider and switches
%     to non-numbered backup frames so "extra / less central" material
%     lives after the main talk without inflating the slide count.
%
% Usage:
%   \documentclass{beamer}
%   % =====================================================================
% finance-beamer.tex — shared Beamer style for the LLM-in-Finance course
% =====================================================================
% Reusable slide style. Each deck \input's this immediately after
% \documentclass{beamer}. It provides:
%   * the Madrid theme + book colour palette (matches the printed book)
%   * two book-style framing blocks:
%       \begin{bigpicture} ... \end{bigpicture}   ("the bigger picture")
%       \begin{underhood}   ... \end{underhood}    ("under the hood")
%     These mirror the book's `context` and `deepdive` tcolorboxes so the
%     same intuition-vs-mechanism scaffold the reader meets in the book
%     reappears on the slides.
%   * a Python listings style for code frames
%   * appendix conventions: \appendixstart drops a divider and switches
%     to non-numbered backup frames so "extra / less central" material
%     lives after the main talk without inflating the slide count.
%
% Usage:
%   \documentclass{beamer}
%   \input{../_style/finance-beamer.tex}
%   \title{...}\subtitle{...}\author{...}\institute{...}\date{\today}
%   \begin{document} ... \end{document}
% =====================================================================

\usetheme{Madrid}
\usecolortheme{whale}
\setbeamertemplate{navigation symbols}{}
\setbeamertemplate{caption}[numbered]
\setbeamercovered{transparent}

% --- Density controls: keep dense, equation-heavy frames on one slide ---
% Slightly smaller body text + tighter lists/blocks. Frame titles and the
% Madrid chrome keep their normal size, so the deck still reads as Madrid,
% just with more vertical room for content.
\setbeamerfont{normal text}{size=\small}
\setbeamertemplate{itemize/enumerate body begin}{\small}
\setbeamertemplate{itemize/enumerate subbody begin}{\footnotesize}
% Tighter block spacing.
\setbeamertemplate{block begin}{%
  \par\vskip2pt%
  \begin{beamercolorbox}[colsep*=.75ex,leftskip=1ex,rightskip=1ex]{block title}%
    \usebeamerfont*{block title}\insertblocktitle%
  \end{beamercolorbox}%
  {\parskip0pt\vskip-1pt}%
  \begin{beamercolorbox}[colsep*=.75ex,vmode,leftskip=1ex,rightskip=1ex]{block body}%
  \vskip1pt}
\setbeamertemplate{block end}{\end{beamercolorbox}\vskip2pt}

\usepackage{amsmath, amssymb}
\usepackage{booktabs}
\usepackage{xcolor}
\usepackage{listings}
\usepackage[skins, breakable]{tcolorbox}

% --- Book colour palette (kept in sync with book/preamble.tex) ---
\definecolor{linkblue}{RGB}{14, 75, 140}
\definecolor{deepdivebg}{RGB}{235, 244, 255}
\definecolor{narrativebg}{RGB}{255, 252, 240}
\definecolor{narrativerule}{RGB}{180, 130, 40}
\definecolor{steelblue}{RGB}{70, 130, 180}
\definecolor{forestgreen}{RGB}{34, 139, 34}
\definecolor{crimson}{RGB}{220, 20, 60}
\definecolor{darkorange}{RGB}{255, 140, 0}
\definecolor{codebg}{RGB}{248, 248, 248}
\definecolor{coderule}{RGB}{210, 210, 210}
\definecolor{keywordblue}{RGB}{0, 100, 180}
\definecolor{commentgray}{RGB}{120, 120, 120}
\definecolor{stringgreen}{RGB}{30, 130, 30}

% Tie the Madrid structural colour to the book's link blue.
\setbeamercolor{structure}{fg=linkblue}
\setbeamercolor{title}{fg=white}
\setbeamercolor{frametitle}{fg=white}

% --- "the bigger picture" : narrative / intuition framing -----------
\newtcolorbox{bigpicture}{%
  enhanced, breakable, boxsep=2pt,
  colback  = narrativebg,
  colframe = narrativerule,
  boxrule  = 0pt, toprule = 1.4pt, bottomrule = 0.4pt,
  leftrule = 0pt, rightrule = 0pt, arc = 0pt,
  left = 6pt, right = 6pt, top = 4pt, bottom = 5pt,
  fonttitle = \footnotesize\scshape\color{narrativerule},
  title = {the bigger picture}, attach title to upper = {\par\smallskip},
}

% --- "under the hood" : the mechanism / technical detail ------------
\newtcolorbox{underhood}{%
  enhanced, breakable, boxsep=2pt,
  colback  = deepdivebg,
  colframe = linkblue,
  boxrule  = 0pt, toprule = 1.4pt, bottomrule = 0.4pt,
  leftrule = 0pt, rightrule = 0pt, arc = 0pt,
  left = 6pt, right = 6pt, top = 4pt, bottom = 5pt,
  fonttitle = \footnotesize\scshape\color{linkblue},
  title = {under the hood}, attach title to upper = {\par\smallskip},
}

% --- Python code style ---------------------------------------------
\lstset{
  basicstyle    = \ttfamily\footnotesize,
  keywordstyle  = \color{keywordblue}\bfseries,
  commentstyle  = \color{commentgray}\itshape,
  stringstyle   = \color{stringgreen},
  showstringspaces = false,
  language      = Python,
  breaklines    = true,
  backgroundcolor = \color{codebg},
  frame         = single,
  rulecolor     = \color{coderule},
  numbers       = none,
  aboveskip     = 4pt, belowskip = 4pt,
}

% --- Appendix conventions ------------------------------------------
% \appendixstart{Title} : begins the "extra material" appendix. Frames
% after it are not counted toward the main deck's frame total, keeping
% the core talk lean while preserving the deeper / optional content.
\newcommand{\appendixstart}[1]{%
  \appendix
  \setbeamertemplate{frametitle continuation}{}%
  \begin{frame}[noframenumbering]
    \begin{center}
      {\Large\scshape\color{linkblue} Appendix}\\[4pt]
      {\large #1}\\[10pt]
      {\footnotesize Extra / optional material — beyond the core lecture.}
    \end{center}
  \end{frame}
}
% Use \begin{frame}[noframenumbering]{...} for each appendix frame.

%   \title{...}\subtitle{...}\author{...}\institute{...}\date{\today}
%   \begin{document} ... \end{document}
% =====================================================================

\usetheme{Madrid}
\usecolortheme{whale}
\setbeamertemplate{navigation symbols}{}
\setbeamertemplate{caption}[numbered]
\setbeamercovered{transparent}

% --- Density controls: keep dense, equation-heavy frames on one slide ---
% Slightly smaller body text + tighter lists/blocks. Frame titles and the
% Madrid chrome keep their normal size, so the deck still reads as Madrid,
% just with more vertical room for content.
\setbeamerfont{normal text}{size=\small}
\setbeamertemplate{itemize/enumerate body begin}{\small}
\setbeamertemplate{itemize/enumerate subbody begin}{\footnotesize}
% Tighter block spacing.
\setbeamertemplate{block begin}{%
  \par\vskip2pt%
  \begin{beamercolorbox}[colsep*=.75ex,leftskip=1ex,rightskip=1ex]{block title}%
    \usebeamerfont*{block title}\insertblocktitle%
  \end{beamercolorbox}%
  {\parskip0pt\vskip-1pt}%
  \begin{beamercolorbox}[colsep*=.75ex,vmode,leftskip=1ex,rightskip=1ex]{block body}%
  \vskip1pt}
\setbeamertemplate{block end}{\end{beamercolorbox}\vskip2pt}

\usepackage{amsmath, amssymb}
\usepackage{booktabs}
\usepackage{xcolor}
\usepackage{listings}
\usepackage[skins, breakable]{tcolorbox}

% --- Book colour palette (kept in sync with book/preamble.tex) ---
\definecolor{linkblue}{RGB}{14, 75, 140}
\definecolor{deepdivebg}{RGB}{235, 244, 255}
\definecolor{narrativebg}{RGB}{255, 252, 240}
\definecolor{narrativerule}{RGB}{180, 130, 40}
\definecolor{steelblue}{RGB}{70, 130, 180}
\definecolor{forestgreen}{RGB}{34, 139, 34}
\definecolor{crimson}{RGB}{220, 20, 60}
\definecolor{darkorange}{RGB}{255, 140, 0}
\definecolor{codebg}{RGB}{248, 248, 248}
\definecolor{coderule}{RGB}{210, 210, 210}
\definecolor{keywordblue}{RGB}{0, 100, 180}
\definecolor{commentgray}{RGB}{120, 120, 120}
\definecolor{stringgreen}{RGB}{30, 130, 30}

% Tie the Madrid structural colour to the book's link blue.
\setbeamercolor{structure}{fg=linkblue}
\setbeamercolor{title}{fg=white}
\setbeamercolor{frametitle}{fg=white}

% --- "the bigger picture" : narrative / intuition framing -----------
\newtcolorbox{bigpicture}{%
  enhanced, breakable, boxsep=2pt,
  colback  = narrativebg,
  colframe = narrativerule,
  boxrule  = 0pt, toprule = 1.4pt, bottomrule = 0.4pt,
  leftrule = 0pt, rightrule = 0pt, arc = 0pt,
  left = 6pt, right = 6pt, top = 4pt, bottom = 5pt,
  fonttitle = \footnotesize\scshape\color{narrativerule},
  title = {the bigger picture}, attach title to upper = {\par\smallskip},
}

% --- "under the hood" : the mechanism / technical detail ------------
\newtcolorbox{underhood}{%
  enhanced, breakable, boxsep=2pt,
  colback  = deepdivebg,
  colframe = linkblue,
  boxrule  = 0pt, toprule = 1.4pt, bottomrule = 0.4pt,
  leftrule = 0pt, rightrule = 0pt, arc = 0pt,
  left = 6pt, right = 6pt, top = 4pt, bottom = 5pt,
  fonttitle = \footnotesize\scshape\color{linkblue},
  title = {under the hood}, attach title to upper = {\par\smallskip},
}

% --- Python code style ---------------------------------------------
\lstset{
  basicstyle    = \ttfamily\footnotesize,
  keywordstyle  = \color{keywordblue}\bfseries,
  commentstyle  = \color{commentgray}\itshape,
  stringstyle   = \color{stringgreen},
  showstringspaces = false,
  language      = Python,
  breaklines    = true,
  backgroundcolor = \color{codebg},
  frame         = single,
  rulecolor     = \color{coderule},
  numbers       = none,
  aboveskip     = 4pt, belowskip = 4pt,
}

% --- Appendix conventions ------------------------------------------
% \appendixstart{Title} : begins the "extra material" appendix. Frames
% after it are not counted toward the main deck's frame total, keeping
% the core talk lean while preserving the deeper / optional content.
\newcommand{\appendixstart}[1]{%
  \appendix
  \setbeamertemplate{frametitle continuation}{}%
  \begin{frame}[noframenumbering]
    \begin{center}
      {\Large\scshape\color{linkblue} Appendix}\\[4pt]
      {\large #1}\\[10pt]
      {\footnotesize Extra / optional material — beyond the core lecture.}
    \end{center}
  \end{frame}
}
% Use \begin{frame}[noframenumbering]{...} for each appendix frame.

%   \title{...}\subtitle{...}\author{...}\institute{...}\date{\today}
%   \begin{document} ... \end{document}
% =====================================================================

\usetheme{Madrid}
\usecolortheme{whale}
\setbeamertemplate{navigation symbols}{}
\setbeamertemplate{caption}[numbered]
\setbeamercovered{transparent}

% --- Density controls: keep dense, equation-heavy frames on one slide ---
% Slightly smaller body text + tighter lists/blocks. Frame titles and the
% Madrid chrome keep their normal size, so the deck still reads as Madrid,
% just with more vertical room for content.
\setbeamerfont{normal text}{size=\small}
\setbeamertemplate{itemize/enumerate body begin}{\small}
\setbeamertemplate{itemize/enumerate subbody begin}{\footnotesize}
% Tighter block spacing.
\setbeamertemplate{block begin}{%
  \par\vskip2pt%
  \begin{beamercolorbox}[colsep*=.75ex,leftskip=1ex,rightskip=1ex]{block title}%
    \usebeamerfont*{block title}\insertblocktitle%
  \end{beamercolorbox}%
  {\parskip0pt\vskip-1pt}%
  \begin{beamercolorbox}[colsep*=.75ex,vmode,leftskip=1ex,rightskip=1ex]{block body}%
  \vskip1pt}
\setbeamertemplate{block end}{\end{beamercolorbox}\vskip2pt}

\usepackage{amsmath, amssymb}
\usepackage{booktabs}
\usepackage{xcolor}
\usepackage{listings}
\usepackage[skins, breakable]{tcolorbox}

% --- Book colour palette (kept in sync with book/preamble.tex) ---
\definecolor{linkblue}{RGB}{14, 75, 140}
\definecolor{deepdivebg}{RGB}{235, 244, 255}
\definecolor{narrativebg}{RGB}{255, 252, 240}
\definecolor{narrativerule}{RGB}{180, 130, 40}
\definecolor{steelblue}{RGB}{70, 130, 180}
\definecolor{forestgreen}{RGB}{34, 139, 34}
\definecolor{crimson}{RGB}{220, 20, 60}
\definecolor{darkorange}{RGB}{255, 140, 0}
\definecolor{codebg}{RGB}{248, 248, 248}
\definecolor{coderule}{RGB}{210, 210, 210}
\definecolor{keywordblue}{RGB}{0, 100, 180}
\definecolor{commentgray}{RGB}{120, 120, 120}
\definecolor{stringgreen}{RGB}{30, 130, 30}

% Tie the Madrid structural colour to the book's link blue.
\setbeamercolor{structure}{fg=linkblue}
\setbeamercolor{title}{fg=white}
\setbeamercolor{frametitle}{fg=white}

% --- "the bigger picture" : narrative / intuition framing -----------
\newtcolorbox{bigpicture}{%
  enhanced, breakable, boxsep=2pt,
  colback  = narrativebg,
  colframe = narrativerule,
  boxrule  = 0pt, toprule = 1.4pt, bottomrule = 0.4pt,
  leftrule = 0pt, rightrule = 0pt, arc = 0pt,
  left = 6pt, right = 6pt, top = 4pt, bottom = 5pt,
  fonttitle = \footnotesize\scshape\color{narrativerule},
  title = {the bigger picture}, attach title to upper = {\par\smallskip},
}

% --- "under the hood" : the mechanism / technical detail ------------
\newtcolorbox{underhood}{%
  enhanced, breakable, boxsep=2pt,
  colback  = deepdivebg,
  colframe = linkblue,
  boxrule  = 0pt, toprule = 1.4pt, bottomrule = 0.4pt,
  leftrule = 0pt, rightrule = 0pt, arc = 0pt,
  left = 6pt, right = 6pt, top = 4pt, bottom = 5pt,
  fonttitle = \footnotesize\scshape\color{linkblue},
  title = {under the hood}, attach title to upper = {\par\smallskip},
}

% --- Python code style ---------------------------------------------
\lstset{
  basicstyle    = \ttfamily\footnotesize,
  keywordstyle  = \color{keywordblue}\bfseries,
  commentstyle  = \color{commentgray}\itshape,
  stringstyle   = \color{stringgreen},
  showstringspaces = false,
  language      = Python,
  breaklines    = true,
  backgroundcolor = \color{codebg},
  frame         = single,
  rulecolor     = \color{coderule},
  numbers       = none,
  aboveskip     = 4pt, belowskip = 4pt,
}

% --- Appendix conventions ------------------------------------------
% \appendixstart{Title} : begins the "extra material" appendix. Frames
% after it are not counted toward the main deck's frame total, keeping
% the core talk lean while preserving the deeper / optional content.
\newcommand{\appendixstart}[1]{%
  \appendix
  \setbeamertemplate{frametitle continuation}{}%
  \begin{frame}[noframenumbering]
    \begin{center}
      {\Large\scshape\color{linkblue} Appendix}\\[4pt]
      {\large #1}\\[10pt]
      {\footnotesize Extra / optional material — beyond the core lecture.}
    \end{center}
  \end{frame}
}
% Use \begin{frame}[noframenumbering]{...} for each appendix frame.


\title{Financial Text Summarization \& Information Extraction}
\subtitle{Lecture 14 --- Turning Filings, Calls, and Reports into Structured, Faithful Facts}
\author{Juan F.\ Imbet}
\institute{Paris Dauphine -- PSL University}
\date{\today}

\begin{document}

\begin{frame}\titlepage\end{frame}

\begin{frame}{Outline}\tableofcontents\end{frame}

% ===========================================================================
\section{The Information Extraction Problem in Finance}
% ===========================================================================

\begin{frame}{Why Extraction and Summarization Matter}
\begin{bigpicture}
Markets run on information, not on its absence. Tens of thousands of filings,
transcripts, and analyst notes are released daily. The bottleneck is the
\emph{cost of reading, interpreting, and routing} --- not the lack of facts.
\end{bigpicture}

\medskip
\textbf{The scale problem.}
\begin{itemize}
  \item A PM cannot read every 10-K from the 6{,}000+ firms in the Russell 3000.
  \item A compliance officer cannot scan every contract in a merger data room.
\end{itemize}

\medskip
\textbf{Three tightly coupled IE sub-tasks:}
\begin{enumerate}
  \item \textbf{Entity recognition} --- who/what the text refers to.
  \item \textbf{Fact / relation extraction} --- how entities relate, what numbers they carry.
  \item \textbf{Summarization} --- condense long documents, faithfully.
\end{enumerate}
\begin{block}{Dependency chain}
NER quality \emph{caps} relation-extraction accuracy; faithful summarization
\emph{requires} accurate entity and fact extraction underneath.
\end{block}
\end{frame}

\begin{frame}{Structured vs.\ Semi- vs.\ Unstructured Text}
\begin{block}{Definition --- the structure spectrum}
A document is \textbf{structured} if every element can be queried programmatically
via a published schema (no NL needed); \textbf{semi-structured} if it mixes
schema-conformant sections with free-form narrative; \textbf{unstructured} if its
content lives in natural language with no imposed schema.
\end{block}

\medskip
\begin{tabular}{lll}
\toprule
\textbf{Tier} & \textbf{Examples} & \textbf{Strategy} \\
\midrule
Structured     & XBRL filings, tick data, options chains & Database query \\
Semi-structured& 10-K / 10-Q, transcripts, proxies       & Tables + NLP \\
Unstructured   & News, social media, analyst notes, contracts & Full NLU \\
\bottomrule
\end{tabular}

\medskip
\begin{alertblock}{Practitioner rule}
Start with structured sources (low noise, low cost). But unstructured text holds
\emph{forward-looking} signal --- guidance, qualitative risk, interpretation ---
absent from the tables. The best pipelines combine both.
\end{alertblock}
\end{frame}

\begin{frame}{Key Entities in Financial Text}
\begin{columns}[T]
\column{0.5\textwidth}
\textbf{Organisations} --- most frequent, most ambiguous. ``Apple'' (firm vs.\
fruit), ``Morgan'' (J.P.\ Morgan vs.\ Morgan Stanley). Resolution =
\emph{entity linking} to LEI / CUSIP / EDGAR CIK.

\medskip
\textbf{Temporal expressions} --- ``next quarter'', ``FY2024'', ``within 30 days of
the trigger''. Must normalise fiscal + calendar terms to dates.

\column{0.5\textwidth}
\textbf{Monetary figures \& metrics} --- \$ amounts, \%, EPS, EBITDA margin, D/E.
Resolve units (M vs.\ B), parenthesised negatives, and the entity/period each
figure describes.

\medskip
\textbf{Legal / contractual} --- covenants in loan agreements and bond indentures
that can trigger default; classify and extract threshold values.
\end{columns}

\medskip
\begin{exampleblock}{Earnings-call excerpt $\to$ extracted records}
``In Q3 2023, \textbf{Microsoft} reported revenue of \textbf{\$56.5 billion}, up
\textbf{13\%} YoY. CEO \textbf{Satya Nadella} noted \textbf{Azure} grew \textbf{29\%}\dots''
\\[2pt]
\textsc{Org}: Microsoft (CIK 789019) $\cdot$ \textsc{Person}: Nadella/CEO $\cdot$
\textsc{Date}: Q3'23 $\to$ 2023-07-01..09-30 $\cdot$ \textsc{Money}: \$56.5B (total rev) $\cdot$
\textsc{Metric}: +13\% rev, +29\% Azure.
\end{exampleblock}
\end{frame}

\begin{frame}{Extraction vs.\ Summarization: A Precise Taxonomy}
\begin{columns}[T]
\column{0.5\textwidth}
\begin{block}{Information Extraction (IE)}
\emph{Schema-driven.} Identify and structure specified information into a
\textbf{predefined schema}. Output = typed records (entities, attributes,
relations). Models are \textbf{discriminative}: classify whether a span fills a slot.
\end{block}
\textit{``What happened to whom, when, by how much?''} $\to$ populates databases.

\column{0.5\textwidth}
\begin{block}{Summarization}
\emph{Content-driven.} Produce a shorter coherent text capturing the salient
content. Output = a \textbf{natural-language string}. Models are
\textbf{generative}; ``salient'' is learned from reference summaries.
\end{block}
\textit{``What should I know without reading it all?''} $\to$ serves readers.
\end{columns}

\medskip
\begin{alertblock}{Why the distinction matters}
The two have \emph{different evaluation criteria and failure modes}. A quant trading
earnings surprises needs IE; an investment committee reviewing a note needs
summarization; a covenant monitor needs IE \emph{then} a human-readable summary.
\end{alertblock}
\end{frame}

% ===========================================================================
\section{Named Entity Recognition in Finance}
% ===========================================================================

\begin{frame}{Financial NER: Beyond the Four CoNLL Classes}
\begin{bigpicture}
NER is the foundational IE task: label each token with its entity type (or O).
The general-domain CoNLL-2003 schema --- \textsc{Person}, \textsc{Location},
\textsc{Organisation}, \textsc{Misc} --- is \emph{inadequate} for finance.
\end{bigpicture}

\medskip
Financial documents add specialised types: \textbf{financial instruments},
\textbf{economic indicators}, \textbf{legal clauses}, \textbf{corporate roles}.
These demand domain-specific datasets and fine-tuned models.

\medskip
\begin{block}{Three finance-specific annotation challenges}
\begin{enumerate}
  \item \textbf{Context-dependent type}: ``Apple'' = org (portfolio) vs.\ product (supply chain).
  \item \textbf{Nested entities}: ``Goldman Sachs Investment Banking Division'' = org $+$ unit.
  \item \textbf{Implicit entities}: ``the issuer'' refers to a firm named earlier.
\end{enumerate}
\end{block}
\end{frame}

\begin{frame}{Financial NER \& RE Datasets}
\small
\begin{tabular}{llll}
\toprule
\textbf{Dataset} & \textbf{Type} & \textbf{Scale / scope} & \textbf{Use case} \\
\midrule
FINER-139 & NER, 139 types  & SEC filings + news      & Fine-grained NER \\
FinNER    & NER, doc-level  & 10-K / 10-Q             & Coreference chains \\
FinRE     & Relations       & (e$_1$, rel, e$_2$) triples & Ownership graphs \\
EDGAR-Corpus & Pretraining  & 6M+ filings, $\sim$6.5B tok & Domain pretraining \\
\bottomrule
\end{tabular}

\medskip
\begin{itemize}
  \item \textbf{FINER-139}: hierarchical taxonomy (Org $\to$ Public Company, ETF,
        Interest Rate\dots). ``Company'' is \emph{not} monolithic --- conflating PE
        firms with listed banks creates systematic graph errors.
  \item \textbf{FinRE} relations: \texttt{Subsidiary-Of}, \texttt{Acquired-By},
        \texttt{Founded-By}, \texttt{Announced-Transaction}.
  \item \textbf{EDGAR-Corpus}: 10-K/10-Q from 1993--2020. Models pretrained on it
        \emph{systematically} beat general-domain models on financial NER.
\end{itemize}
\end{frame}

\begin{frame}{Fine-Tuning a Transformer NER Head}
\begin{underhood}
Token classification on a pretrained encoder. For tokens
$\mathbf{x}=(x_1,\dots,x_T)$ with contextual states $\mathbf{h}_i \in \mathbb{R}^d$:
\[
P(y_i \mid \mathbf{x}) = \mathrm{softmax}(\mathbf{W}\mathbf{h}_i + \mathbf{b}),
\quad \mathbf{W}\in\mathbb{R}^{C\times d}.
\]
Labels follow \textbf{BIO}: $\{B\text{-}e, I\text{-}e, O\}$. Train end-to-end on
token cross-entropy $\mathcal{L}=-\sum_t \log p(\hat y_t = y_t)$.
\end{underhood}

\medskip
\textbf{Encoder choice matters.} FinBERT $>$ vanilla BERT; \textbf{SEC-BERT}
(pretrained only on EDGAR) is SOTA on SEC-domain NER --- it has seen
``Regulation S-K Item 1A'' thousands of times.

\medskip
\textbf{Three finance-specific recipe tweaks beyond the standard:}
\begin{enumerate}
  \item \textbf{Segmentation}: docs exceed 512 tokens $\to$ sliding windows
        ($\geq$64-token overlap) or Longformer.
  \item \textbf{Label imbalance}: O is 85--95\% of tokens $\to$ focal / weighted loss.
  \item \textbf{Post-processing}: a CRF + Viterbi layer enforces valid BIO transitions.
\end{enumerate}
\end{frame}

\begin{frame}{Example: SEC-BERT $\to$ FINER-139}
\begin{exampleblock}{Production-grade recipe}
\begin{enumerate}
  \item Download FINER-139 splits ($\sim$900 annotated training sentences).
  \item Init from SEC-BERT checkpoint (MLM-pretrained on millions of SEC sentences).
  \item Linear head of dimension $C = 2\times 139 + 1 = 279$ (BIO for 139 types $+$ O).
  \item AdamW ($\beta_1{=}0.9,\beta_2{=}0.999$, wd $0.01$), linear warmup 10\%,
        peak LR $2\times10^{-5}$, 10 epochs.
\end{enumerate}
\end{exampleblock}

\medskip
\textbf{Result:} $\sim$20 min on one A100; macro-F$_1$ \emph{substantially} exceeds
the FinBERT baseline on rare types ($<$50 training examples) --- where domain
pretraining pays off most.
\end{frame}

\begin{frame}{Relation Extraction: Entities $\to$ Facts}
\begin{bigpicture}
Entities acquire financial meaning through relationships: Microsoft
\emph{acquired} Activision; Buffett \emph{is CEO of} Berkshire; a loan is
\emph{collateralised by} specific assets. These feed knowledge graphs,
earnings attribution, and corporate-hierarchy databases.
\end{bigpicture}

\begin{underhood}
Joint entity-and-relation models (e.g.\ DyGIE++) score ordered span pairs:
\[
s(e_i, r, e_j) = f\bigl(\mathbf{h}_{e_i}, \mathbf{h}_{e_j}, \mathbf{r}\bigr),
\]
with span reps $\mathbf{h}_{e}$ (mean/endpoint of encoder states), learned
relation embedding $\mathbf{r}\in\mathbb{R}^{d_r}$, and a bilinear/FFN scorer $f$;
predict $r \in \mathcal{R}\cup\{\text{None}\}$.
\end{underhood}

\medskip
\textbf{Taxonomy:} \textsc{Subsidiary-Of}, \textsc{Acquired-By},
\textsc{Employed-At}, \textsc{Board-Member-Of}, \textsc{Issued-By},
\textsc{Secured-By}, \textsc{Reported-Value}.
\\[2pt]
\textbf{Hard case:} \emph{cross-sentence} RE $\gg$ within-sentence. Document-level
models (DocRED-style) that encode full context win.
\end{frame}

% ===========================================================================
\section{Financial Document Summarization}
% ===========================================================================

\begin{frame}{Extractive vs.\ Abstractive --- The Faithfulness Tradeoff}
\begin{columns}[T]
\column{0.5\textwidth}
\begin{block}{Extractive}
Select a subset of sentences/spans \textbf{verbatim} from the source.
Faithfulness \textbf{guaranteed by construction} --- every sentence was written by
the original author.
\end{block}
\column{0.5\textwidth}
\begin{block}{Abstractive}
Generate \textbf{new text} that paraphrases and synthesises. More concise and
fluent --- but susceptible to \textbf{hallucination}: plausible-sounding but
factually wrong statements.
\end{block}
\end{columns}

\medskip
\begin{alertblock}{Why hallucination is consequential in finance}
Summarising ``revenue of \$57B'' when the truth was \$56.5B is small in \% but can
be \emph{material misinformation} in a decision process. SOTA abstractive models
hallucinate factual errors in \textbf{25--30\%} of CNN/DailyMail summaries; untuned
financial models fare no better.
\end{alertblock}

\medskip
\textbf{Recommendation:} prefer extractive when faithfulness is paramount
(esp.\ numbers); use abstractive when fluency dominates. \textbf{Hybrid} (extract,
then rephrase) is a useful middle ground.
\end{frame}

\begin{frame}{Earnings \& 10-K Summarization}
A typical 10-K runs \textbf{50{,}000--150{,}000 words}: business, risk factors,
legal proceedings, MD\&A, financial statements, notes.

\medskip
\textbf{The MD\&A} is the densest section for an equity investor --- management's
interpretation and forward guidance. Its tone, readability, and length carry
measurable signal about future earnings and returns (Loughran--McDonald).
Tone-preservation must be an \emph{explicit} constraint in any faithful prompt.

\medskip
\begin{block}{Three-stage pipeline}
\begin{enumerate}
  \item \textbf{Section extraction \& routing} --- parse TOC, route each section to
        a sub-model (table extractor vs.\ text summariser).
  \item \textbf{Chunk-level summarization} --- 512-tok chunks (64 overlap) for BERT;
        whole MD\&A in one pass for 8K--128K decoders.
  \item \textbf{Hierarchical aggregation} --- second-pass or retrieval-based merge of
        chunk summaries into a document summary.
\end{enumerate}
\end{block}
\end{frame}

\begin{frame}[fragile]{Example: Constrained Prompt for Earnings}
\begin{exampleblock}{Numerical-hallucination-resistant prompt}
\begin{verbatim}
You are a senior equity analyst. Summarise the following earnings
release in no more than 150 words, covering: (1) revenue and EPS
vs. consensus, (2) the single most important forward guidance
statement, and (3) one key risk mentioned by management.
Quote all numbers exactly as stated.

[EARNINGS RELEASE TEXT]
\end{verbatim}
\end{exampleblock}

\medskip
\textbf{The key line:} ``Quote all numbers exactly as stated'' is a practical
prompt-engineering technique that reduces numerical hallucination.

\medskip
\textbf{Illustrative finding} (figures vary with corpus/model): the constrained
prompt cut number errors by $\sim$40\% vs.\ an unconstrained paraphrase prompt;
$\sim$8\% of cases still showed paraphrase errors on nuanced guidance ---
motivating the \emph{factual-consistency} checks later.
\end{frame}

\begin{frame}{Analyst Notes \& Long-Document Challenges}
\begin{columns}[T]
\column{0.52\textwidth}
\textbf{Analyst-note condensation.} Reports run 10--40 pp with a predictable
macro-structure (thesis in para 1, target price standardised). Template-aware
parsing exploits this.
\\[3pt]
\textbf{Change detection / delta summarization}: a PM wants \emph{what changed},
not what they already know --- requires report history + reasoning over diffs.

\column{0.48\textwidth}
\textbf{Earnings transcripts}: 15{,}000--30{,}000 words; uneven density
(prepared remarks dense; Q\&A mixes signal + boilerplate). \textbf{ECTSum}: 2{,}425
transcript--summary pairs, 3--6 abstractive bullets each.
\\[3pt]
\textbf{S-1 filings}: a tech S-1 can exceed \textbf{300{,}000 words}.
\end{columns}

\medskip
\textbf{Common thread:} no single context window fits these documents, so the
architecture --- not just the prompt --- must handle length. Three strategies
follow on the next slide.
\end{frame}

\begin{frame}{Long-Document Strategies}
\begin{underhood}
Three strategies for extreme-length documents:
\begin{itemize}
  \item \textbf{Two-stage hierarchical}: each section $\to$ 200--400 words, then a
        second-pass model over the concatenated section summaries.
  \item \textbf{Retrieval-augmented summarization}: dense-retrieve top-$k$ passages
        for a query, summarise only those (summarization as RAG).
  \item \textbf{Long-context models}: Longformer; 128K--1M-token models
        (e.g.\ Gemini 1.5 Pro) process an entire S-1 in one pass --- cost-bound.
\end{itemize}
\end{underhood}
\end{frame}

% ===========================================================================
\section{Table and Numerical Data Extraction}
% ===========================================================================

\begin{frame}{SEC XBRL: Structured Extraction Without NLP}
\begin{bigpicture}
Tables are 2-D grids where each cell's meaning depends on its row header (what
measure) \emph{and} column header (what period/entity). Reliable table extraction
underpins database population and consistency checking.
\end{bigpicture}

\medskip
\textbf{Inline XBRL} (mandated for U.S.\ domestic filers since 2020): every figure
tagged with a U.S.-GAAP concept (e.g.\ \texttt{us-gaap:Revenues}), a context
(entity + period), and units --- machine-readable \emph{and} browser-readable.

\smallskip
\begin{block}{Prefer XBRL whenever available}
The EDGAR XBRL API returns a figure with full provenance for a (CIK, concept,
year) query. Accuracy is guaranteed by the filer and its auditors --- \emph{no NLP}.
\end{block}

\smallskip
\begin{alertblock}{Coverage gap}
Non-GAAP measures (adjusted EBITDA, organic growth, book-to-bill) and many
footnotes are \emph{not} XBRL-tagged. Combine XBRL for core figures with
LLM/NLP extraction for the surrounding narrative.
\end{alertblock}
\end{frame}

\begin{frame}{LLMs for Table Understanding}
\textbf{When XBRL is unavailable} (international filings, private docs, pre-XBRL
legacy): tables must be parsed from HTML/PDF/text --- merged cells, multi-row
headers, embedded footnote refs, non-standard layouts.

\medskip
\begin{underhood}
\textbf{FinQA} benchmarks multi-step numerical reasoning over tables:
\begin{itemize}
  \item Baseline cell-retrieval models: $\sim$50\% exact-match.
  \item Fine-tuned models that \emph{generate a Python-like calc program}: $>$70\%.
\end{itemize}
\end{underhood}

\medskip
\textbf{Why generation beats retrieval:} financial table questions chain several
operations (subtract prior year, divide, express as \%). A model that emits an
explicit calculation program externalises that reasoning and routes the arithmetic
to a deterministic interpreter --- continued next slide.
\end{frame}

\begin{frame}{LLMs for Table Understanding --- Pipeline}
\textbf{Four-stage LLM pipeline} for parsing tables when XBRL is unavailable:
\begin{enumerate}
  \item \textbf{Layout parsing} (LayoutLM / PDF parser) $\to$ cells with row/col coords.
  \item \textbf{Header disambiguation} $\to$ LLM labels each cell with (row-hdr, col-hdr).
  \item \textbf{Value normalisation} $\to$ units, parenthesised negatives, \% vs.\ absolute.
  \item \textbf{Schema mapping} $\to$ map to GAAP taxonomy concepts.
\end{enumerate}

\medskip
\begin{alertblock}{Arithmetic reliability (ReAct principle)}
LLMs are unreliable calculators. \textbf{Route all arithmetic to a code interpreter}:
the LLM decides \emph{what} to compute; Python \emph{actually} computes it. This
separation dramatically cuts numerical errors.
\end{alertblock}
\end{frame}

\begin{frame}{Consistency Checking: Numbers vs.\ Prose}
The same revenue figure appears in the tables, the press release, and the MD\&A.
\textbf{Consistency checking} verifies they agree --- catching (a) \emph{extraction
errors} and (b) \emph{genuine disclosure inconsistencies} (possibly material).

\medskip
\begin{block}{Three-step procedure}
\begin{enumerate}
  \item \textbf{Multi-source extraction} --- pull every occurrence of each metric
        (revenue, op.\ income, EPS) from tables, press release, MD\&A, footnotes.
  \item \textbf{Normalisation} --- common unit \& sign. ``\$56.5B'' $\to$
        \$56{,}500M; all sources should agree.
  \item \textbf{Cross-reference} --- flag any pair disagreeing beyond tolerance
        (e.g.\ 0.1\% for rounding); report with source attribution.
\end{enumerate}
\end{block}

\medskip
\textbf{Prose vs.\ table fact-checking.} ``Net income grew 25\% YoY'' can be
verified against extracted current/prior figures --- a powerful quality gate for
both document pipelines \emph{and} LLM-generated summaries, where hallucinated
numbers are a known failure mode.
\end{frame}

% ===========================================================================
\section{Evaluation and Benchmarks}
% ===========================================================================

\begin{frame}{ROUGE --- Recall-Oriented Overlap}
\begin{bigpicture}
\textbf{ROUGE} is the most widely used automatic summarization metric. Intuitively:
how much of the \emph{reference} summary's wording does the candidate reproduce?
A recall-oriented overlap score.
\end{bigpicture}

\begin{underhood}
ROUGE-$N$ = $n$-gram recall of candidate $c$ against references $\mathcal{R}$
($\mathcal{G}_N(s)$ = $n$-grams of $s$):
\[
\text{ROUGE-}N=\frac{\sum_{s\in\mathcal{R}}\sum_{g\in\mathcal{G}_N(s)}
\mathrm{Count}_{\text{match}}(g,c)}
{\sum_{s\in\mathcal{R}}\sum_{g\in\mathcal{G}_N(s)}\mathrm{Count}(g)}.
\]
ROUGE-L uses the longest common subsequence instead of exact $n$-grams.
\end{underhood}
\end{frame}

\begin{frame}{ROUGE --- Why Lexical Overlap Is Not Enough}
ROUGE rewards \emph{surface} $n$-gram overlap with the reference. That makes it
cheap and reproducible, but blind to meaning.

\medskip
\begin{alertblock}{Limitation: lexical, not semantic}
``net income declined 12\%'' $\to$ ``profit fell 12\%'' deserves full credit but
scores \textbf{zero} ROUGE-1 on the substituted words. Conversely, copying
irrelevant source sentences can score \emph{high} ROUGE without being informative.
\end{alertblock}

\medskip
\textbf{Takeaway:} treat ROUGE as a coarse recall proxy, not a quality verdict ---
it motivates the semantic and factual metrics that follow.
\end{frame}

\begin{frame}{BERTScore --- Semantic Matching}
\begin{bigpicture}
\textbf{BERTScore} fixes ROUGE's lexical-overlap blind spot: candidate and
reference tokens count as equivalent when their \emph{meanings} match, not only
their surface forms.
\end{bigpicture}

\begin{underhood}
Cosine similarity of contextual embeddings for candidate token $\hat r_j$ and
reference token $r_i$:
\[
\mathrm{sim}(r_i,\hat r_j)=
\frac{\mathbf{h}_{r_i}^\top \mathbf{h}_{\hat r_j}}
{\|\mathbf{h}_{r_i}\|\,\|\mathbf{h}_{\hat r_j}\|}.
\]
Precision/recall/F$_1$ via greedy token matching in both directions. With a
\emph{financial} encoder (FinBERT), it captures domain-specific similarity better.
\end{underhood}
\end{frame}

\begin{frame}{BERTScore --- What Semantic Matching Still Misses}
BERTScore closes ROUGE's paraphrase gap, but similarity in embedding space is
still not the same as being \emph{factually correct}.

\medskip
\begin{alertblock}{Neither ROUGE nor BERTScore measures factual correctness}
A summary inventing ``revenue grew 15\%'' (true: 13\%) can still score high if the
surrounding language matches the reference. And outputs can be inconsistent
\emph{across runs} --- reproducibility itself is a QA concern for production IE.
\end{alertblock}

\medskip
\textbf{Takeaway:} similarity metrics rank fluency and coverage, not truth ---
hence the factual-consistency checks on the next slide.
\end{frame}

\begin{frame}{Factual Consistency \& Human Evaluation}
\begin{block}{Factual-consistency metrics}
\textbf{FactCC} trains an NLI classifier (entailment / contradiction / neutral) to
predict whether each candidate claim is \emph{entailed} by the source. For finance,
strengthen by \textbf{grounding numerical claims against extracted XBRL}: any claim
that cannot be verified against structured data is flagged.
\end{block}

\medskip
\textbf{Human evaluation --- the gold standard --- by \emph{domain professionals}}
(a crowdworker cannot judge GAAP vs.\ non-GAAP revenue). Assess four dimensions:
\begin{itemize}
  \item \textbf{Factual accuracy} (most critical; mark each claim verified / unverifiable / incorrect)
  \item \textbf{Completeness} --- all material disclosures covered?
  \item \textbf{Conciseness} --- no redundancy?
  \item \textbf{Fluency / readability}
\end{itemize}

\medskip
\textbf{Report agreement} (Cohen's $\kappa$ / Krippendorff's $\alpha$): typically
\emph{low} on completeness (analysts value different facts), \emph{high} on factual
accuracy (errors are objective) $\Rightarrow$ accuracy is the tractable auto target.
\end{frame}

\begin{frame}{Inventory of Financial NLP Datasets}
\small
\begin{tabular}{llll}
\toprule
\textbf{Dataset} & \textbf{Document type} & \textbf{Scale} & \textbf{Primary task} \\
\midrule
ECTSum       & Earnings transcripts & 2{,}425 pairs ($\sim$7k words ea.) & Long-doc summarization \\
EDGAR-Corpus & SEC filings          & 6M+ filings        & Domain pretraining \\
FINER-139    & SEC filings + news   & 139 entity types   & Fine-grained NER \\
FinSBD       & Financial text       & annotated          & Sentence boundary detection \\
FinQA        & Tables + text        & multi-step Q\&A    & Numerical reasoning \\
\bottomrule
\end{tabular}

\medskip
\begin{alertblock}{No single benchmark covers the domain}
A model that tops ECTSum (summarization) may fail on FinQA (table reasoning).
Read ``SOTA on financial NLP'' \emph{with the specific benchmark in mind}. Match
the dataset to your document type and task --- and construct a \emph{held-out} test
set from documents outside any public benchmark before deploying.
\end{alertblock}
\end{frame}

\begin{frame}{Lecture 14 --- Five Takeaways}
\begin{enumerate}
  \item \textbf{Financial text is heterogeneous} --- structured (query) to
        unstructured (full NLU). Production pipelines combine both, anchored on
        structured sources.
  \item \textbf{Financial NER needs domain data \& models} --- general-domain NER
        is inadequate for the 139-class taxonomy. Pretrain on EDGAR, fine-tune on
        FINER-139.
  \item \textbf{Abstractive summarization carries hallucination risk} --- in
        finance, hallucinated numbers can be \emph{misinformation}. Pair with
        factual-consistency checking against structured data.
  \item \textbf{Tables \& numerical consistency are underappreciated} --- XBRL is
        free and accurate; LLMs supplement for non-GAAP \& legacy; route arithmetic
        to code.
  \item \textbf{Match the metric to the stakes} --- ROUGE/BERTScore measure
        similarity, not facts. Add factual-consistency checks \emph{and} expert
        human evaluation.
\end{enumerate}
\begin{block}{Next}
Practical 14: an end-to-end SEC 10-K pipeline --- section routing, extractive
summarization, ROUGE scoring, and consistency checking.
\end{block}
\end{frame}

% ===========================================================================
% APPENDIX
% ===========================================================================
\appendixstart{Extended Material --- Annotation, RE Internals, Further Reading}

\begin{frame}[noframenumbering]{Appendix: Annotation Challenges in Finance}
General-domain NER guidelines do not anticipate these:
\begin{itemize}
  \item \textbf{Surface-string ambiguity}: ``Apple'' as organisation (portfolio)
        vs.\ product (supply-chain discussion) --- the \emph{same string}, different type.
  \item \textbf{Nested entities}: ``Goldman Sachs Investment Banking Division'' =
        org (Goldman Sachs) $+$ business unit, simultaneously.
  \item \textbf{Implicit entities}: ``the issuer'' in a bond prospectus refers to a
        previously named company without restating it.
\end{itemize}
\medskip
These motivate \emph{domain-specific annotation guidelines and adjudication
processes} --- not just relabelled CoNLL guidelines.
\end{frame}

\begin{frame}[noframenumbering]{Appendix: Why Cross-Sentence RE Is Hard}
\begin{itemize}
  \item Within-sentence RE: both entities co-occur $\to$ local context suffices.
  \item Cross-sentence RE: entities in different sentences $\to$ requires
        \emph{document-level} encoding and coreference resolution.
  \item Benchmarks: DocRED and its financial adaptations; document-level models
        that encode the whole context (not sentence-by-sentence) achieve the gains.
\end{itemize}
\medskip
\textbf{Why it matters in finance:} ownership chains and ``the issuer''-style
implicit references routinely span paragraphs in prospectuses and 10-Ks.
\end{frame}

\begin{frame}[noframenumbering]{Appendix: Retrieval-Augmented Summarization}
For documents that exceed any context window (S-1 $>$300k words), turn
summarization into \textbf{retrieval-augmented generation}:
\begin{enumerate}
  \item Index the document's passages with a dense retriever.
  \item For a query (``What are the principal risk factors?''), retrieve top-$k$.
  \item Summarise \emph{only} the retrieved passages.
\end{enumerate}
\medskip
\textbf{Tradeoff vs.\ hierarchical summarization:} RAG is query-targeted and cheap
but may miss globally salient content not matching the query; hierarchical
two-stage covers everything but costs a full pass over the document.
\end{frame}

\begin{frame}[noframenumbering]{Appendix: Further Reading}
\begin{itemize}
  \item \textbf{NER foundations}: Lample et al.\ (2016); Devlin et al.\ (2019) survey.
        Financial: FINER-139 (Loukas 2022); English/Greek FINER (Shah 2023).
  \item \textbf{Summarization}: ECTSum (Mukherjee 2022) for transcripts;
        EDGAR-Corpus (Loukas 2021) for domain-adaptive pretraining.
  \item \textbf{Evaluation}: ROUGE (Lin 2004); BERTScore (Zhang 2020); FactCC
        (Kryscinski 2020); faithfulness analysis (Maynez 2020).
  \item \textbf{Table reasoning}: FinQA (Chen 2021); broad survey FinBen (Xie 2024).
\end{itemize}
\medskip
\textbf{Cross-chapter links:} XBRL + table parsing here are prerequisites for the
DCF pipeline in the Business Valuation chapter; apply the hallucination-mitigation
framework (LLM Limitations chapter) to any decision-critical abstractive summary.
\end{frame}

\end{document}
